% 08_conclusion.tex: Section 7.
\section{Conclusion}\label{sec:conclusion}

I asked whether the delay of an arrival at AIBD, defined as \blockon{} later than STA by more than 15 minutes, can be predicted from information that is known before the aircraft arrives. For this purpose, I built a chronological evaluation on the 2025 ledger of \num{10863} flights, restricted the models to seven schedule variables, and tested them on the last month of the year.

The main finding is that the schedule carries some, but limited, information about delays. On the December test set, the logistic regression reaches a macro F1 of 0.668 and a ROC AUC of 0.749, and gradient boosting reaches 0.644 and 0.739. Both models are clearly better than the majority baseline, which has a macro F1 of 0.315, and both are clearly above the no-skill level in the precision-recall analysis. The logistic regression is ahead of gradient boosting in all three robustness variants, but the difference is small (0.024 macro F1, with an interval from 0.005 to 0.042), and the two models were practically tied on the validation month. I therefore do not conclude that the more flexible model has an advantage on these data.

The error analysis and the threshold analysis show where the models are weak. They miss many delays in the morning and for several airlines, and at the threshold of 0.5 they flag fewer flights than were delayed in December. Moving the threshold to 0.45 improves the macro F1 of both models in this month, but I report it only as a diagnosis, and a threshold for operational use would need the costs of the two kinds of error.

What the results support is a first screening tool: a model of this kind could rank the scheduled arrivals of the next days by their risk of delay, so that planners can look at the flights with the highest scores first. What they do not support is a reliable warning for single flights, because 37\% to 40\% of the delays occurred without a flag, and because the evidence comes from one test month at one airport (\cref{sec:limitations}).

Three steps would strengthen the evidence. The first is a rolling evaluation over several test months, so that the variation between months is measured and not only the variation between days. The second is to add information that is available before arrival, for example the status of the inbound flight and the weather forecast, which would address the main limit of a schedule-only model. The third is to choose the threshold from the operational costs of false alarms and missed delays, and to check the calibration of the scores before they are used for decisions.
